\documentclass[11pt]{article}
\usepackage[a4paper,margin=1in]{geometry}
\usepackage{amsmath,amssymb,mathtools,bm}
\usepackage{enumitem,booktabs,array}
\usepackage{tcolorbox}
\usepackage{hyperref}
\usepackage[protrusion=true,expansion=false]{microtype}
\hypersetup{colorlinks=true,linkcolor=blue,urlcolor=blue,citecolor=blue}
\setlist[itemize]{topsep=2pt,itemsep=2pt}
\setlist[enumerate]{topsep=2pt,itemsep=2pt}
\newtcolorbox{keybox}{colback=blue!3!white,colframe=blue!40!black,boxrule=0.4pt,arc=1mm}
\newtcolorbox{alertbox}{colback=red!3!white,colframe=red!40!black,boxrule=0.4pt,arc=1mm}
\newtcolorbox{answerbox}{colback=green!3!white,colframe=green!40!black,boxrule=0.4pt,arc=1mm}
\newtcolorbox{drillbox}{colback=yellow!5!white,colframe=orange!60!black,boxrule=0.4pt,arc=1mm}
\newcommand{\EV}{\mathbb{E}}
\newcommand{\Var}{\mathrm{Var}}
\newcommand{\Cov}{\mathrm{Cov}}
\newcommand{\blank}[1]{\underline{\hspace{#1}}}

\title{W09-01: Gormley, Matsa \& Milbourn (2013, JAE) \\
\large ``CEO Compensation and Corporate Risk: Evidence from a Natural Experiment'' --- Active Recall Workbook}
\author{Banking \& Risk Management --- Exam Prep}
\date{\today}
\begin{document}\maketitle

\begin{keybox}
\textbf{Paper in one sentence.} After unexpected EPA regulatory actions added a large \emph{left-tail liability risk} to certain industries (mid-1980s to early 1990s), boards at affected firms reduced CEO compensation \emph{vega} (risk-taking incentive) by roughly 20\% relative to unaffected firms --- showing that boards \emph{do} adjust risk-taking incentives in response to exogenous increases in downside risk.

\textbf{Placement.} Canonical evidence that executive compensation convexity (vega, delta) is \emph{endogenous} to the firm's risk environment. Pairs with Bakke et al.\ (W09-02) on FAS 123R as a different comp-shock.
\end{keybox}

\section*{Round 1: Complete Walk-Through}

\subsection*{1.1 The shock}
Between 1986 and 1990, EPA actions (CERCLA/Superfund, expanded enforcement) significantly raised liability risk for firms in certain chemical-using industries (SIC-based classification). The extra risk is heavily \emph{left-tailed}: rare but catastrophic litigation/cleanup outcomes.

\subsection*{1.2 Compensation measures}
\begin{itemize}
\item \textbf{Delta.} Sensitivity of CEO portfolio value to 1\% change in stock price. Rewards performance.
\item \textbf{Vega.} Sensitivity of CEO portfolio value to 1\% change in stock volatility. Rewards risk-taking.
\end{itemize}
ExecuComp data.

\subsection*{1.3 Identification}
DiD across affected vs.\ unaffected SIC codes, pre vs.\ post regulatory event:
\[
\text{Vega}_{it}=\alpha_i+\delta_t+\beta\cdot(\text{Affected}_i\cdot\text{Post}_t)+X_{it}'\gamma+\varepsilon_{it}.
\]
Expected $\beta<0$: boards cut vega to reduce risk-taking incentives in now-riskier firms.

\subsection*{1.4 Headline results}
\begin{itemize}
\item Vega at affected firms falls by $\sim 20\%$ post-shock relative to control firms.
\item Delta moves little --- boards still care about performance alignment.
\item Result is driven by \emph{new} grants, not vesting of legacy grants (consistent with active board re-optimisation).
\item Placebos on pseudo-events find no effect.
\end{itemize}

\subsection*{1.5 Interpretation}
\begin{itemize}
\item Endogeneity of comp-structure is first-order: regressions that treat vega/delta as exogenous will be misspecified.
\item Supports a view that comp-design is an active \emph{risk-management} tool of the board.
\item Links neatly to Smith-Stulz (1985) / Stulz (1996) on value of risk management under downside costs.
\end{itemize}

\subsection*{1.6 Caveats}
\begin{alertbox}
\begin{itemize}
\item Identification relies on the orthogonality of EPA-affected-industry assignment to comp trends.
\item DiD pre-trend check is necessary.
\item CEO compensation is a slow-moving contract; short-run effects may understate long-run.
\end{itemize}
\end{alertbox}

\section*{Round 2: Level 1 Fill-in}
\begin{enumerate}
\item Shock: \blank{1cm} regulation raising \blank{2cm}-tail liability risk.
\item Vega change at affected firms: $\sim$\blank{1cm}\% \blank{1cm}.
\item Delta change: approximately \blank{1cm}.
\item Design: \blank{1cm}-in-\blank{1cm} across industries.
\item Data source for comp: \blank{2cm}.
\end{enumerate}

\section*{Round 3: Level 2 Prompts}
\begin{enumerate}
\item Define delta and vega formally and explain what each incentivises.
\item Write the DiD specification and state the identifying assumption.
\item Why do boards cut vega but leave delta alone?
\item What does GMM imply for papers that treat comp vega as exogenous?
\end{enumerate}

\section*{Round 4: Minimal Scaffolding}
From a blank page: (i) shock, (ii) vega/delta definitions, (iii) DiD, (iv) headline magnitude, (v) implication for endogeneity of comp.

\section*{Round 5: Blank Page}
Essay: is CEO compensation a risk-management tool? GMM evidence.

\vspace{6pt}
\begin{drillbox}
\textbf{Speed Drill.} (1) Regulatory source. (2) Tail targeted. (3) Vega change direction and magnitude. (4) Delta change. (5) Data source.
\end{drillbox}

\section*{Quick Reference \& Answer Key}
\begin{answerbox}
\begin{itemize}
\item \textbf{Shock.} EPA/CERCLA raising left-tail liability risk, late 1980s.
\item \textbf{Design.} DiD affected vs.\ unaffected SIC industries.
\item \textbf{Vega.} $\sim\!20\%$ fall in affected firms.
\item \textbf{Delta.} Unchanged.
\end{itemize}
\end{answerbox}

\begin{keybox}
\textbf{30-second exam pitch.} Gormley, Matsa \& Milbourn (2013) use a natural experiment --- the late-1980s EPA/CERCLA expansion that abruptly raised left-tail liability risk for firms in certain chemical-using industries --- to show that boards actively retune CEO compensation in response to risk. Vega (sensitivity of CEO wealth to volatility) falls by about 20\% in affected firms relative to unaffected peers after the shock, while delta (performance sensitivity) is unchanged. The result comes from new grants, consistent with deliberate rebalancing. Implication: CEO compensation structure is endogenous to the firm's risk environment, and regressions that treat comp as exogenous are misspecified. Theoretically, the result is in the Smith--Stulz tradition: boards use comp-convexity to align the CEO's risk-taking incentives with the value of corporate risk management.
\end{keybox}

\end{document}
